\section{Contratos y limites}

\subsection{Descripcion}

Este documento delimita que puede saber y hacer cada capa del sistema. Su objetivo es evitar acoplamiento accidental entre dominio, frontend, backend, datos e integraciones.

\subsection{Alcance}

Aplica a los limites entre repositorios y entre capas internas principales.

\subsection{Definiciones}

\begin{itemize}
  \item \textbf{Contrato}: acuerdo explicito de entrada, salida y comportamiento esperado.
  \item \textbf{Limite}: frontera de responsabilidad entre componentes.
\end{itemize}

\subsection{Reglas}

Contratos entre repositorios:

\begin{itemize}
  \item El frontend consume datos del backend por HTTP y contratos aplicacionales.
  \item El frontend no depende directamente de la base de datos ni de la API externa de resultados.
  \item El backend usa la base de datos como sistema de registro y persistencia.
  \item El backend encapsula integraciones externas y expone al frontend una vista de dominio estable.
\end{itemize}

Limites internos esperados del backend:

\begin{itemize}
  \item `auth`: autenticacion y sesion.
  \item `users`: identidad y estado de cuenta.
  \item `payments`: recepcion y validacion funcional de comprobantes.
  \item `matches`: calendario, kickoff y estados de partido.
  \item `picks`: captura y lock de picks.
  \item `results`: recepcion de resultados oficiales y overrides.
  \item `scoring`: calculo de puntos y desempates.
  \item `leaderboard`: vistas de ranking.
  \item `admin`: acciones administrativas.
  \item `audit`: trazabilidad.
  \item `notifications`: correos y avisos del sistema.
\end{itemize}

Contratos de documentacion:

\begin{itemize}
  \item Las reglas funcionales se definen primero en \texttt{01\_Producto\_y\_Reglas}.
  \item Los datos conceptuales se definen primero en \texttt{05\_Datos}.
  \item Backend y frontend consumen estas definiciones, no las reemplazan.
\end{itemize}

\subsection{Casos borde}

\begin{itemize}
  \item Si una integracion externa entrega un modelo incompatible con el dominio, se adapta en el limite de integracion, no en el frontend.
  \item Si una regla de negocio obliga a recalculo, el contrato del dominio prevalece sobre cualquier cache o vista derivada.
\end{itemize}

\subsection{Invariantes}

\begin{itemize}
  \item Ninguna capa consume directamente internals de otra sin contrato definido.
  \item Ninguna integracion externa define la logica oficial del juego por si sola.
\end{itemize}

\subsection{Preguntas abiertas}

\begin{itemize}
  \item Confirmar si notifications vivira como modulo independiente o como servicio transversal.
\end{itemize}

\subsection{Decisiones pendientes}

\begin{itemize}
  \item `Decision pendiente` `No bloqueante`: confirmar el limite exacto entre `admin` y los modulos de dominio cuando se cierre la arquitectura backend.
\end{itemize}
